% Generated draft: requires mathematical and source review.
\documentclass[11pt]{article}
\usepackage[margin=25mm]{geometry}
\usepackage{fontspec}
\setmainfont{Latin Modern Roman}
\usepackage{amsmath,amssymb,mathtools,unicode-math}
\setmathfont{Latin Modern Math}
\usepackage{enumitem,needspace,xurl,hyperref}
\setlength{\parindent}{0pt}
\setlength{\parskip}{5pt}
\setlength{\emergencystretch}{3em}
\newcommand{\N}{\mathbb{N}}
\newcommand{\Z}{\mathbb{Z}}
\newcommand{\Q}{\mathbb{Q}}
\newcommand{\R}{\mathbb{R}}
\newcommand{\C}{\mathbb{C}}
\newcommand{\F}{\mathbb{F}}
\newcommand{\A}{\mathbb{A}}
\newcommand{\PP}{\mathbb{P}}
\newcommand{\E}{\mathbb{E}}
\newcommand{\eps}{\varepsilon}
\newcommand{\dd}{\,\mathrm{d}}
\begin{document}
\section*{UnsolvedMath ID: 10900054}
\textbf{Source problem code:} catalogue entry
\subsection{Definitions and mathematical statement}
Let \(M\) be a finite-volume hyperbolic \(3\)-manifold. In the complete quotient interpretation write \(M=\mathbb H^3/\Gamma\), where \(\Gamma\) is a lattice in the appropriate hyperbolic isometry group. Call \(M\) arithmetic when \(\Gamma\) is commensurable with an arithmetic lattice, up to conjugacy and the usual orientation convention.

\emph{Question (McMullen).} If \(M\) contains infinitely many immersed totally geodesic surfaces, must \(M\) be arithmetic? An intended geometric formulation requires infinitely many genuinely distinct surface images, not just infinitely many parametrizations or finite covers of one immersed surface. With an admissible class of surfaces fixed, it asks whether
\[
 \#\{\text{distinct immersed totally geodesic surface images in }M\}=\infty
 \quad\Longrightarrow\quad M\text{ is arithmetic}.
\]
The term ``surface'' must include the intended completeness and area conditions: arbitrary local geodesic patches would not express the arithmeticity problem. The source excerpt does not state whether the surfaces are required to be closed or may be noncompact of finite area, nor how distinctness is counted. Its short formulation is therefore preserved with those necessary conventions unresolved.

The supplied registry labels the problem solved. Its summary credits Margulis and Mohammadi with an arithmeticity theorem for finite-volume hyperbolic \(3\)-manifolds containing infinitely many immersed, closed totally geodesic surfaces, and reports a later broader generalization by Bader--Fisher--Miller--Stover. These are supplied literature claims, not independently verified theorems here; in particular the summary's added word ``closed'' must be compared with the broader source wording.
\subsection{Short English statement}
Does an infinite collection of distinct immersed totally geodesic surfaces force a finite-volume hyperbolic \(3\)-manifold to be arithmetic? The registry reports a solution, but the source's surface completeness, finite-area, closedness, and distinctness conventions need to be matched to the cited result.
\subsection{Sources}
\begin{itemize}
\item Local UnsolvedMath registry, permanent ID 10900054, source code catalogue entry; attribution McMullen, subsection 7.6, ``Totally geodesic surfaces and arithmeticity.''
\item Delp, Hoffoss and Manning, \emph{Problems In Groups, Geometry, and Three-Manifolds} (2015), subsection 7.6, PDF page 6. Supplied URLs: \url{https://arxiv.org/pdf/1512.04620#page=6} and \url{https://arxiv.org/abs/1512.04620}; recorded access 2026-07-29; extraction source-tex checked against rendered PDF text. No source-side resolution is recorded.
\item The research summary cites Margulis and Mohammadi, ``Arithmeticity of hyperbolic 3-manifolds containing infinitely many totally geodesic surfaces,'' GAFA 2021, \texttt{arXiv:1907.05815}; and Bader--Fisher--Miller--Stover, ``Arithmeticity, superrigidity, and totally geodesic subspaces,'' \texttt{arXiv:2004.10255}. It describes the latter as a higher-rank and more-general-subspace generalization. Exact hypotheses and bibliographic identifications have not been independently checked.
\item Supplied status solved and SOLVED-IN-LITERATURE. The generic background instead records the 2015 open presentation and requires a complete modern primary-literature audit; this is compatible with later progress in principle but does not itself verify the claimed match. Status evidence and redistribution terms of the publicly accessible author-submitted TeX and PDF remain NEEDS\_REVIEW.
\end{itemize}
\end{document}
